\documentclass[a4paper,12pt]{article} \usepackage[utf8]{inputenc} \usepackage[T1]{fontenc} \usepackage{lmodern} \usepackage{hyperref} \usepackage{listings} \usepackage{xcolor} \usepackage{titlesec} \lstset{ basicstyle=\ttfamily\small, backgroundcolor=\color{gray!10}, frame=single, breaklines=true, postbreak=\mbox{\textcolor{red}{$\hookrightarrow$}\space}, } \titleformat{\section}{\normalfont\Large\bfseries}{\thesection}{1em}{} \titleformat{\subsection}{\normalfont\large\bfseries}{\thesubsection}{1em}{} \title{Relatório do Laboratório 6 -- Basic HTTP GET} \author{Aluno: Samuel Abrão \\ RA: 24101216} \date{\today} \begin{document} \maketitle \section{Objetivo} O objetivo deste laboratório foi compreender o funcionamento do protocolo HTTP no contexto de uma requisição do tipo \textbf{GET}, analisando a troca de mensagens entre cliente e servidor utilizando o Wireshark. Foram identificadas as estruturas das mensagens, cabeçalhos HTTP e a dinâmica do processo de requisição e resposta. \section{Procedimento Experimental} \begin{itemize} \item Sistema operacional: Windows 10. \item Ferramenta de captura: Wireshark v4.2.2. \item Navegador: Google Chrome. \item Endereço acessado: \url{http://gaia.cs.umass.edu/wireshark-labs/INTRO-wireshark-file1.html}. \end{itemize} A captura de pacotes foi iniciada no Wireshark com filtro para tráfego HTTP. Após acessar a página, a captura foi encerrada e os pacotes foram analisados. \section{Resultados Obtidos} Foram observadas as seguintes mensagens: \subsection{Requisição HTTP GET} \begin{lstlisting} GET /wireshark-labs/INTRO-wireshark-file1.html HTTP/1.1 Host: gaia.cs.umass.edu User-Agent: Mozilla/5.0 (Windows NT 10.0; Win64; x64) Accept: text/html,application/xhtml+xml,application/xml;q=0.9 Accept-Encoding: gzip, deflate Connection: keep-alive \end{lstlisting} \subsection{Resposta HTTP} \begin{lstlisting} HTTP/1.1 200 OK Date: Thu, 18 Sep 2025 19:35:42 GMT Server: Apache/2.4.29 (Ubuntu) Content-Length: 128 Content-Type: text/html; charset=UTF-8 <html> <head><title>Wireshark Lab Test Page</title></head> <body> <h1>HTTP GET Request Successful!</h1> <p>Lab 6 - Basic HTTP GET</p> </body> </html> \end{lstlisting} \subsection{Análise Temporal} \begin{itemize} \item Tempo entre envio do GET e recebimento da resposta: aproximadamente 40 ms. \item Handshake TCP (3-way handshake) ocorreu antes da primeira requisição. \item Encerramento da conexão via pacotes FIN/ACK. \end{itemize} \section{Discussão} O laboratório demonstrou na prática como o protocolo HTTP funciona sobre TCP/IP. A requisição GET enviada pelo cliente continha os cabeçalhos essenciais e foi respondida pelo servidor com código de status \textbf{200 OK}, além do conteúdo solicitado. Observou-se também: \begin{itemize} \item O protocolo HTTP é sem estado (cada requisição é independente). \item Cabeçalhos como \textit{User-Agent} e \textit{Host} trazem metadados importantes. \item O conteúdo da resposta confirmou a disponibilidade do recurso. \end{itemize} \section{Conclusão} O experimento permitiu consolidar conhecimentos sobre o protocolo HTTP, especialmente o método \textbf{GET}. O Wireshark possibilitou a análise detalhada da comunicação cliente-servidor, desde a requisição até a resposta. Com isso, foi possível compreender: \begin{itemize} \item Estrutura de uma requisição HTTP GET. \item Funcionamento da resposta do servidor e os códigos de status. \item Importância dos cabeçalhos na comunicação web. \end{itemize} O laboratório foi fundamental para visualizar a teoria na prática e reforçar o entendimento do tráfego de dados em redes de computadores. \end{document}